% Lookup in a 2-way set-associative cache with 4 sets (8 lines, 16-byte
% blocks, 32-bit addresses): the index selects a set, both tags are compared,
% and each comparison is ANDed with the valid bit of its line.
\begin{tikzpicture}[x=1cm,y=1cm, font=\footnotesize\sffamily, >={Stealth[length=1.8mm]},
  gate/.style={draw, fill=ln-box, minimum width=0.7cm, minimum height=0.5cm, inner sep=1pt, font=\sffamily\scriptsize},
  cmp/.style={draw, circle, fill=white, inner sep=1pt, minimum size=0.5cm}]
  % --- address fields
  %% The Japanese index label filled its box edge to edge, so the Japanese
  %% tag/index boundary sits further left (the index line at x=5.4 and the
  %% tag line at x=1.0 stay inside their boxes).
  \iflangja{%
    \draw[fill=ln-accent!15] (0.6,0) rectangle (3.6,0.55);
    \node at (2.1,0.275) {タグ\ (31\,\dots\,6)};
    \draw[fill=ln-box] (3.6,0) rectangle (6.4,0.55);
    \node at (5.0,0.275) {インデックス\ (5\,\dots\,4)};
  }{%
    \draw[fill=ln-accent!15] (0.6,0) rectangle (4.0,0.55);
    \node at (2.3,0.275) {tag\ (31\,\dots\,6)};
    \draw[fill=ln-box] (4.0,0) rectangle (6.4,0.55);
    \node at (5.2,0.275) {index\ (5\,\dots\,4)};
  }
  \draw[fill=white] (6.4,0) rectangle (8.8,0.55);
  \node at (7.6,0.275) {\iflangja{オフセット}{offset}\ (3\,\dots\,0)};
  % --- two ways, four sets; set 0 selected
  \foreach \x in {1.6,5.8} {
    \fill[ln-accent!15] (\x,-1.1) rectangle (\x+3.4,-1.55);
    \node[font=\sffamily\scriptsize, text=ln-muted, anchor=south] at (\x+0.2,-1.1) {V};
    \node[font=\sffamily\scriptsize, text=ln-muted, anchor=south] at (\x+1.1,-1.1) {\iflangja{タグ}{tag}};
    \node[font=\sffamily\scriptsize, text=ln-muted, anchor=south] at (\x+2.6,-1.1) {\iflangja{データ}{data}};
    \foreach \i in {0,...,3} {
      \draw (\x,-1.1-0.45*\i) rectangle (\x+0.4,-1.55-0.45*\i);
      \draw (\x+0.4,-1.1-0.45*\i) rectangle (\x+1.8,-1.55-0.45*\i);
      \draw (\x+1.8,-1.1-0.45*\i) rectangle (\x+3.4,-1.55-0.45*\i);
    }
  }
  \node[font=\sffamily\scriptsize, text=ln-muted, anchor=north] at (4.2,-2.95) {\iflangja{ウェイ 0}{way 0}};
  \node[font=\sffamily\scriptsize, text=ln-muted, anchor=north] at (8.4,-2.95) {\iflangja{ウェイ 1}{way 1}};
  \foreach \i in {0,...,3} {
    \node[font=\sffamily\scriptsize, text=ln-muted, anchor=west] at (9.25,-1.325-0.45*\i) {\iflangja{セット}{set}\ \i};
  }
  % --- index selects a set in both ways
  \draw (5.4,0) -- (5.4,-1.325);
  \fill (5.4,-1.325) circle (1.2pt);
  \draw[->] (5.4,-1.325) -- (5.0,-1.325);
  \draw[->] (5.4,-1.325) -- (5.8,-1.325);
  % --- per way: tag comparator, AND with the valid bit
  \foreach \x/\w in {1.6/0,5.8/1} {
    \node[cmp] (c\w) at (\x+1.2,-3.6) {$=$};
    \node[gate] (a\w) at (\x+0.4,-4.35) {AND};
    \draw[->] (\x+1.2,-2.9) -- (c\w.north);
    \draw[->] (\x+0.2,-2.9) -- (\x+0.2,-4.1);
    \draw[->] (c\w.west) -| (\x+0.6,-4.1);
  }
  % --- tag of the address to both comparators
  \draw (1.0,0) -- (1.0,-5.5) -- (7.0,-5.5);
  \draw[->] (7.0,-5.5) -- (c1.south);
  \fill (2.8,-5.5) circle (1.2pt);
  \draw[->] (2.8,-5.5) -- (c0.south);
  % --- either way hits
  \node[gate] (or) at (4.5,-4.35) {OR};
  \draw[->] (a0.east) -- (or.west);
  \draw[->] (a1.west) -- (or.east);
  \draw[->] (or.south) -- (4.5,-4.95) node[below] {\iflangja{ヒット}{hit}};
\end{tikzpicture}
